Sei $\ell$ die Länge des Brecheisens und $a$ der Abstand des Drehpunkts von der Spitze. Der Hebelarm der Kraft $F_\mathrm{L}$ auf die Spitze ist $a$, der Hebelarm der Hand wird ebenfalls vom Drehpunkt aus gemessen: $\ell-a$. Die beiden Kräfte drehen das Brecheisen in entgegengesetzte Richtungen, im Gleichgewicht sind ihre Drehmomente gleich gross:
\begin{align}
 F_\mathrm{H}\cdot(\ell-a) &= F_\mathrm{L}\cdot a
\end{align}
Also
\begin{align}
 F_\mathrm{H} &= \FHF \\
   &= \frac{\FL\cdot\aS}{\lB-\aS} \\
   &= \FH \approx \result{\FHP{0}{2}}
\end{align}
Ein langer Hebelarm macht aus einer kleinen Kraft eine grosse.
